% !PRINTED 536
\weberpage{536}

From such a basis form one obtains all the integers of the field
$\Omega$ if one puts integral rational numbers for the variables.

The fundamental form $\tau$ is the root of an irreducible equation
of degree $n$
\[
F(t) = N(t - \tau) = 0,
\]
in which the coefficients of the powers of $t$ are integral rational (and
homogeneous) functions of the variables $t_1, t_2, \ldots, t_n$.

\begin{center}
\S. 147.

\gesperrt{The absolute norms of the functionals.}
\end{center}

The elements $\alpha_1, \alpha_2, \ldots, \alpha_n$ of a basis of the
functional $\mu$ can, as integers in $\Omega$, be expressed linearly
and with integral coefficients through a basis $\omega_1, \omega_2,
\ldots, \omega_n$ of $\rho$ in the form
\weq{1}{
(\alpha_1, \alpha_2, \ldots, \alpha_n) = A(\omega_1, \omega_2, \ldots,
\omega_n),
}
where $A$ denotes a linear substitution with integral rational
coefficients. A special case of this is offered by the formulae
(3) of the preceding section.

A basis form $\lambda$ of $\mu$ we shall represent as follows:
\weq{2}{
\lambda = \sum^{\nu} \alpha_\nu t_\nu,
}
and we write the substitution (1) more fully:
\weq{3}{
\alpha_s = \sum^{\nu} a_{s,\nu}\,\omega_\nu,
}
where the $t_\nu$ are variables, the $a_{s,\nu}$ are the substitution
coefficients, and the summation letter $\nu$ runs from 1 to $n$.

Since the products $\alpha_s\,\omega_r$ are all divisible by $\mu$, they
can, according to the meaning of the basis, be represented in the form
\weq{4}{
\alpha_\nu\,\omega_r = \sum^{s} g^{(s)}_{\nu,r}\,\alpha_s,
}
where the $g^{(s)}_{\nu,r}$ are integral rational numbers. From this one
then obtains, by (2),
\weq{5}{
\lambda\,\omega_r = \sum^{s} \alpha_s\,t_{s,r},
}
if
\weq{6}{
t_{s,r} = \sum^{\nu} g^{(s)}_{\nu,r}\,t_\nu
}
are integral rational linear forms.

If one again substitutes the expressions (3) in (5), it follows that
\weq{7}{
\lambda\,\omega_r = \sum^{\nu} \omega_\nu \sum^{s} a_{s,\nu}\,t_{s,r}.
}
